% BEGIN THERMAL_R03_SYNC_13
\section{Conclusiones}
\label{sec:ii-conclusiones}

La construcción obtiene las secciones de acción de Planck y su cociente
de retorno, evalúa el parámetro de Barbero--Immirzi desde la incidencia
orientada y produce los ángulos CKM con su fase y amplitudes unitarias.
Determina asimismo el transductor de Boltzmann para una sección térmica
positiva y el coeficiente gravitatorio común de la realización R1--R8,
cuya especialización circular conserva la coincidencia de radios.
La inscripción determina \(r_N=5759/23040\); la composición longitudinal
da \(L=r_N\alpha_{\HMT}^{16}L_*\); la métrica radial y la acción
fijan \(G_{\rm ret}=c_{\rm int}^3L^2/\hbar_{\rm ret}\) bajo la
identificación expresa de radio reducido. Las identidades resultantes enlazan las constantes
con sus funciones: energía por frecuencia, resolución elemental de área,
información de mezcla, conversión térmica e invariancia geométrica.

La relación central consiste en la conservación de la escala geométrica.
El producto \(G_a\hbar_a\) conserva la escala areal mientras la acción
y la gravitación se transforman recíprocamente; la incidencia ponderada
por \(\gamma_{\HMT}\) determina la resolución de esa área; y
\(k_B\Theta_{{\rm clk},a}\log3=h_a/T_{108}\) relaciona la misma acción
con la energía de una decisión trítica. La observación conjunta del área
y del Hamiltoniano modular, con \(\Delta_{\rm mod}\ne0\) y
normalización conocida, recupera las ocupaciones sectoriales que
una sola lectura agregada no distingue. La inversión directa de
\((N,D)\) conserva su validez sin esa condición. Cada identidad conserva las
condiciones de su realización dimensional, geométrica o espectral.

\subsection{Constantes de Planck, Barbero--Immirzi y mezcla CKM}

La escala de acción contiene dos determinaciones diferenciadas. El
cociente local del registro tiene dieciséis hojas; la norma multiplicativa
y la autoescala fijan \(\operatorname{ord}_{10}(\varphi\alpha^{16})=-34\). Los caracteres regionales fijan las secciones orientadas
y el cociente de retorno. La conversión entre acción angular y acción
por vuelta mantiene la identidad \(h=2\pi\hbar\). La coordenatización angular
transporta los representantes en grados a radianes conservando tanto
los argumentos de los canales como los coeficientes de los funcionales.
Esta precisión permite seguir la dependencia de la acción sin una
inversión retrospectiva de sus coordenadas.

La transición hexádico--octádica aporta el contenido propio del
funcional de Barbero--Immirzi. Los espacios de incidencia fijan los
cardinales; la vuelta orientada fija las multiplicidades; la evaluación
de los canales determina el escalar. Su utilización posterior en el
operador de área y en la acción gravitatoria hace visible el significado
de la estructura que produce el valor. El parámetro ocupa así una
función demostrativa en la relación entre incidencia, orientación y
geometría.

La construcción CKM determina la mezcla mediante un lector explícito. El
lector racional transporta las coordenadas angulares anteriores a tres
ángulos de mezcla y fija su fase; su inversa reconstruye las coordenadas
de partida. La factorización unitaria explicita las amplitudes y el
invariante de Jarlskog conserva la información de violación CP bajo los
cambios de fase admisibles. Su comparación cuantitativa se realiza con
la convención angular y la referencia física declaradas. La construcción
leptónica se presenta en sus marcos espectrales propios, conservando la
diferencia entre los dos transportes de especie.

Los resultados de mezcla quedan expresados numéricamente por
\begin{align*}
 (\theta_{12},\theta_{23},\theta_{13})
   &=(13.003313589509,\ 2.398056157332,\ 0.213977382244)^\circ,\\
 \delta_{\rm CP}&=65.676173123554^\circ,\qquad
 J_{\rm CKM}=3.1189723326632974\times10^{-5}.
\end{align*}
La tabla~\ref{tab:ii-ckm-comparacion} compara sus senos, fase e
invariante con PDG 2025 en la misma parametrización: las cinco
desviaciones marginales normalizadas satisfacen $|z|<0.029$.
El resultado conecta las coordenadas generadas con amplitudes unitarias
y con una medida de fase que permanece invariante bajo reparametrización.

La realización angular en la incidencia de Witt establece un transporte
isométrico explícito, normalizado polarmente, con inversa sobre su imagen.
La igualdad métrica del corolario~\ref{cor:ii-witt-composed-reader}
precisa su función: la composición con el lector angular conserva la
métrica inducida, distingue coordenadas angulares diferentes y permite
reconstruir las tres coordenadas que alimentan las reglas sectoriales.
La comparación entre el registro y la incidencia se realiza así mediante
una aplicación fiel. La representación de determinante conserva la
orientación; los cambios de calibre transforman conjuntamente los datos,
según \ref{prop:ii-witt-gauge-covariance}; la realización compleja conserva
la fase y el orden de los factores unitarios de mezcla. Quedan diferenciadas
la selección de las operaciones sectoriales, su evaluación y el transporte
de su resultado, sin identificar el transporte isométrico posterior con la selección de
las reglas sectoriales que evalúa.

\Needspace{11\baselineskip}
\subsection{Constante de Boltzmann y energía por decisión trítica}

La acción y el período producen
\(\epsilon_a=h_a/T_{108}=\hbar_a/t_P\); la información de la terna
uniforme fija la energía por nat \(\epsilon_a/\log3\).
Fijada una sección térmica positiva \(\Theta_{{\rm clk},a}\),
existe un único transductor lineal positivo
\(k_B:\mathcal L_\Theta\to\mathcal L_E\) que la transforma en esa
energía. Su expresión en bases positivas conserva la ley
\(b'=(s_\Theta/s_E)b\). Así quedan determinados tanto la función
dimensional de Boltzmann como las condiciones que individualizan su
coordenada: sección térmica, bases y, cuando se reivindica selección
interna, el lector que las fija.

El mismo transductor actúa sobre las dos orientaciones cuando las
temperaturas del ciclo tienen el cociente de las secciones de acción,
según \eqref{eq:ii-kb-aut-orientaciones}. Asigna a la información
su magnitud entrópica mediante \(S=k_Bs\), sin identificar una distribución
de continuaciones con una bipartición cuántica ni una actualización
reversible con un borrado físico. La energía, la información y la
entropía conservan así sus dominios al relacionarse con el área y
con la celda gravitatoria.

La construcción de capacidad de \ref{sec:ii-ter27-antecedentes} conserva el calendario
de capacidad con sus vacancias, las dos graduaciones del portador y la
realización energética positiva de su memoria. El espectro decimal, el
origen de capacidad y las multiplicidades determinan la intensidad
\(b_*=1.380649\times10^{-23}\) en la asignación de grados explicitada.
Las corrientes fijan el desfase en su realización conectada; el transporte
centrado conserva el lector con su referencia.

El teorema~\ref{thm:ii-ter27-transducer} añade una consecuencia
termométrica precisa. La referencia fija y el estado de ensayo producen
un incremento entrópico \(b_*s\); su lectura conjunta con la energía
condicionada determina \(\Theta_P=1/(b_*\beta)\) y un único
transductor positivo de coeficiente \(b_*\). El
corolario~\ref{cor:ii-ter27-variational} demuestra además el mínimo
único de energía libre. La referencia fija y el par de funcionales
constituyen las condiciones de esa realización. En el portador graduado,
el Hamiltoniano cronológico de capacidad da exactamente
\(\rho=q^N/Z_N(q)\); su enlace con área y horizonte conserva las
escalas distintas y la identidad finita de entropía relativa.

La sección~\ref{sec:ii-aditiva-instrumento} incorpora la caracterización
aditiva de ese peso y la realización conservativa de sus marcas.
La reunión de alternativas disjuntas, su medida recibida y la unidad de
referencia fijan la contribución $b_*s$; el condicionamiento recupera
la energía media del mismo ensayo. El instrumento conserva el complemento
espectral, la energía del sistema y el transporte trítico con cociente y acarreo.
Una referencia compartida mantiene un único factor $b_*$ al componer
sistemas y conserva sus correlaciones. Estas pruebas explicitan las
operaciones de la realización térmica, además de evaluar su coeficiente.

\subsection{Gravitación y conservación de la escala areal}

La construcción gravitatoria conserva también una información tensorial
que la sección escalar de acoplamiento no agota. El descenso de incidencia
de \ref{sec:ii-pentadica}, bajo sus condiciones explícitas, produce el
portador de cinco componentes. El entrelazador \(\Gamma_5\) realiza ese
portador como tensores simétricos sin traza; su base distingue los pesos
\(\pm2,\pm1,0\). La reducción transversal conserva dos de ellos mediante
un proyector de rango dos. Esta secuencia transmite simultáneamente el
acoplamiento común \(G_a\), la orientación de sus cinco componentes y
la observación reducida, conforme a
\eqref{ii:pent:eq:gravity-12-5-5-2}. La interpretación física conserva la
dinámica y las restricciones especificadas en ese desarrollo.

La inscripción, el archivo completo y las composiciones longitudinal
y métrica determinan primero el lector radial. La energía y el radio
actúan sobre el mismo carácter positivo; su cancelación demuestra que
el coeficiente de acoplamiento es común a todos los modos del dominio
declarado. Las reglas R4 y R5 figuran como composiciones incorporadas,
y R8 mantiene la identificación física del radio reducido. El reloj
\(t_0=\pi_{\HMT}L/(54c_{\rm int})\) expresa posteriormente la
misma longitud en el calendario circular.

El transporte unitario conserva el coeficiente y los registros. La
eliminación conjunta de frontera y memoria conserva la proporción
entre radio y energía, el producto radial y el residuo recuperable;
en el resolvente completo se transporta también el parámetro espectral,
según \eqref{g26:eq:resolvente-proporcional-es}. Esta conservación
abarca la memoria estática y la dinámica dentro de sus dominios.
La reciprocidad de las secciones de acción conserva el
producto de gravitación y acción y, con él, la escala de área. La
realización térmica enlaza la energía de la duración elemental con la
información trítica. En la celda de horizonte especificada, energía,
temperatura, entropía y media acción satisfacen las identidades
demostradas. El artículo sitúa así la relación entre información y
gravedad dentro de una composición matemática explícita.

\subsection{Información de frontera y entropía de entrelazamiento}

La lectura conjunta de área e información sectorial proporciona una
expresión concreta de la genealogía de frontera. El área agrega las
ocupaciones, mientras el observable ponderado conserva su distribución
entre los dos sectores. Su recuperación se demuestra mediante la
inversión del mapa de incidencia. La entropía de entrelazamiento añade
el estado reducido y la bipartición que le otorgan significado. El
transporte de estos datos permite distinguir conservación espectral,
conservación de orientación y conservación de la información accesible
en una reducción.

La formulación sectorial obtiene además las medias del área y de la
información modular a partir de la función de partición finita. Sus
derivadas se expresan mediante varianzas y covarianzas de ocupación.
La inversión unitaria de las ocupaciones conserva la entropía y
transforma el área en su complemento. Queda precisado así un fenómeno
central de la lectura genealógica: una cantidad conservada puede
coexistir con una transformación estructural que otro observable
detecta. El transductor de Boltzmann conserva su función dimensional
en esa relación entre información, energía y entropía.

El transporte de los modos colectivos explicita cómo la incidencia
hexádico--octádica alcanza las etiquetas areales. La primera ley
modular relaciona, a estado de referencia fijo, las variaciones de
entropía y de ocupación ponderada. Para cambios finitos, la identidad
incluye la entropía relativa entre los estados; conservar ese término
precisa el contenido informacional que distingue una variación
infinitesimal de una transformación finita. La recuperación conjunta
de las ocupaciones mediante área e incidencia sitúa esta relación
directamente sobre los dos observables de frontera.

\subsection{Reciprocidad y continuación excepcional}

La elipse de acción y la dualidad T prolongan el argumento mediante
invariantes de transporte. El intercambio de semiejes conserva el
determinante de la forma; la orientación simpléctica distingue
realizaciones con la misma energía; el intercambio momento--enrollamiento
conserva el operador espectral bajo la identificación de dominios
correspondiente. La realización excepcional mantiene la incidencia
hasta Moonshine y permite formular la reciprocidad modular en su
espacio propio.

La realización de covarianza de la
sección~\ref{sec:ii-area-covarianza} precisa el contenido operatorio de
esa reciprocidad. En la normalización de dos desviaciones estándar,
el área \(h\) fija
\(\det V_\eta=\hbar^2/4\), mientras que la anisotropía angular determina
\(\Delta P/\Delta Q=e^{-\eta}\). El estado gaussiano construido y sus
operadores canónicos realizan exactamente esas dos relaciones; la
rotación simpléctica intercambia sus dispersiones conservando el
conmutador. De este modo quedan relacionados el determinante de acción,
la geometría de la elipse y la saturación de incertidumbre en una
representación explícita. La distinción entre la energía del contorno,
\(\hbar\omega\), y la energía media del estado,
\(\hbar\omega/2\), preserva la normalización de ambas lecturas.

La prolongación a teoría M requiere especificar la compactificación,
las superficies de mundo y las representaciones físicas correspondientes.
La relación partícula--antipartícula mediante orientación, adelanto y
retardo constituye una propuesta de extensión, distinta de las
identidades de covarianza y reciprocidad demostradas aquí.

La sección~\ref{delta22:ii:transporte} demuestra la conservación
exacta de la acción cuadrática bajo transporte entre elipses, deriva
su término de conexión hamiltoniana y establece la condición de
positividad del generador. La versión conformemente simpléctica
conserva la acción normalizada y explicita la forma transportada.

\subsection{Evaluación del registro y conservación de la observación}

La evaluación incidencial de
\eqref{eq:registro-composicion-incidencial} hace explícita la operación
que convierte un libro de visitas y trazas orientadas en veinticinco
coordenadas. Cada terna de recuentos se representa mediante sus residuos
decimales, mientras sus cocientes y procedencias permanecen en el libro.
La proposición~\ref{prop:registro-36-residuos} determina cuándo dos
colecciones de recuentos tienen la misma imagen. Su contenido permite
precisar qué información utiliza el valor obtenido y qué información
adicional conserva la historia que lo produjo. La colección de incidencias
y las trazas que alimentan el evaluador conservan sus condiciones de
construcción; la invertibilidad posterior del registro tiene una función
complementaria, la recuperación exacta de esa representación.

La conservación de la observación adquiere una formulación dinámica en
\ref{subsec:ii-memoria-resolvente}. La identidad finita retiene el estado
terminal de la historia; su prolongación controla la contribución de
los incrementos posteriores mediante una cota de cola. Al eliminar
un sector auxiliar, las identidades
\eqref{mem:schur-fuente}--\eqref{mem:schur-lector} transportan también
la fuente y la aplicación de lectura. El contenido de una trayectoria
puede así expresarse en menos componentes conservando los efectos de
sus excursiones y retornos. La transitividad de
\eqref{mem:schur-transitivo} asegura la coherencia de esa operación
cuando la eliminación se realiza en varias etapas. Esta es una
consecuencia precisa de tratar la memoria como parte operativa del estado.

% BEGIN THERMAL_R03_SYNC_13B
% Formalización reunida de resultados ya expuestos en el cuerpo.
\subsection{Definiciones estructurales y resultados cuantitativos}
\label{sec:ii-conclusiones-cuantitativas}

La definición por estructura conserva el objeto anterior a la lectura
escalar. Para las constantes de acción, ese objeto es una sección orientada
con una norma multiplicativa sobre el cociente local y un carácter regional.
Para Barbero--Immirzi, es
la evaluación de una incidencia orientada. Para la mezcla, es un transporte
entre coordenadas y amplitudes. Para Boltzmann, es una aplicación entre
rectas de magnitudes. Para la gravitación, es una sección recíproca de la
acción cuya magnitud procede de la longitud inscrita y de la lectura
radial especificadas en R1--R8. Estas definiciones
determinan operaciones diferentes dentro de la misma genealogía.

El cuadro siguiente fija, para cada magnitud, el objeto que la define,
la operación que determina su lectura y la función física desarrollada
en el cuerpo. Las pruebas de reconstrucción y transporte establecen
cómo se conserva esa información al reutilizar el resultado. En
particular, \ref{prop:registro-36-residuos} caracteriza la representación
del registro, \ref{mem:prop-schur} conserva la observación bajo reducción
y \ref{cor:ii-witt-composed-reader} conserva las coordenadas angulares
en la realización incidencial. Estas tres operaciones tienen dominios
propios y funciones demostrativas distintas dentro de la misma secuencia.

\begingroup
\setlength{\LTpre}{0pt}
\setlength{\LTpost}{2pt}
\begin{longtable}{@{}p{.13\linewidth}p{.415\linewidth}p{.395\linewidth}@{}}
\caption{Objeto estructural, lectura producida y significado físico de
los resultados del artículo.}\label{tab:ii-conclusiones-estructura}\\
\toprule
Magnitud & Definición estructural y resultado & Lectura física y desarrollo interno\\
\midrule
\endfirsthead
\toprule
Magnitud & Definición estructural y resultado & Lectura física y desarrollo interno\\
\midrule
\endhead
\bottomrule
\endfoot
Planck & Sección de la recta de acción:
$\hbar_{\rm pre}=H_5\,10^{-34}\mathcal U_S$,
$\hbar_{\rm ret}=\eta_{\rm ret}\,10^{-34}\mathcal U_S$.
El cociente local tiene dieciséis hojas y
$\operatorname{ord}_{10}(\varphi\alpha^{16})=-34$.
& La sección relaciona frecuencia y energía; la vuelta completa determina
$h_a=2\pi\hbar_a$. Desarrollo en \ref{sec:action}; coordenadas y comparación
en \ref{tab:ii-planck-comparacion}.\\[5pt]
Barbero--Immirzi & Evaluación orientada de los canales y de la incidencia:
$\gamma_{\HMT}=\pi AC^*/180+12\Delta S_{90}-\Delta S_{120}$.
La evaluación impresa es $0.2740076726944592747\ldots$.
& El coeficiente interviene en la resolución elemental de área y en
la realización Cartan--Holst. Las definiciones de
\ref{sec:barbero} y \ref{sec:area} conservan estos dos usos;
\eqref{eq:holst} especifica el segundo.\\[5pt]
Mezcla CKM & Imagen del vector angular mediante las tres aplicaciones
sectoriales; fase $\delta_{\deg}=9A_{\deg}$.
Los ángulos y amplitudes se evalúan después de construir el lector.
& El cuarteto de Jarlskog expresa la información de fase invariante
bajo cambios diagonales de base. Las reglas, su evaluación y la
realización unitaria se desarrollan en
\ref{sec:ii-propagacion-angular}.\\[5pt]
Boltzmann & Transducción positiva
$k_B:\mathcal L_\Theta\to\mathcal L_E$ con
$k_B\Theta_{{\rm clk},a}\log3=h_a/T_{108}$.
El estado y la medida de las continuaciones determinan el contenido
informacional sobre el que actúa.
& $S=k_Bs$ realiza dimensionalmente la entropía del estado reducido.
Las bases de energía y temperatura y sus transformaciones se conservan
en \ref{sec:ii-boltzmann}; la caracterización aditiva y el instrumento
se demuestran en \ref{sec:ii-aditiva-instrumento}; la comparación está en
\ref{tab:ii-kb-g-cartas}.\\[5pt]
Realización térmica marcada & Referencia fija $\eta$, con multiplicidades
$(d_0,d_1,d_2)=(1,380,649)$ en grados $23,26,29$ y traza
$b_*=1.380649\times10^{-23}$.
Las lecturas $\mathcal S_P=b_*s$ y $\mathcal E_P=\operatorname{Tr}(\rho H)$
determinan $d\mathcal S_P/d\mathcal E_P=b_*\beta$ y
$\Theta_P=(b_*\beta)^{-1}$.
& Con referencia fija, bases positivas transportadas y Hamiltoniano no
escalar, \ref{thm:ii-ter27-transducer} determina el transductor.
El corolario~\ref{cor:ii-ter27-variational} prueba el mínimo único.
La composición conserva $q^N/Z_N(q)$, el origen energético y los
transportes reloj--horizonte de \eqref{eq:ii-ter27-intertwiner}.\\[5pt]
Gravitación & Sección recíproca
$G_a=c_{\rm int}^{3}L^2/\hbar_a$,
con $L=(5759/23040)\alpha_{\HMT}^{16}L_*$.
El reloj $t_0=\pi_{\HMT}L/(54c_{\rm int})$ conserva la expresión
circular $G_a=\vartheta_{108}^{\circ}c_{\rm int}^{5}t_0^2/\hbar_a$,
con $\vartheta_{108}^{\circ}=(54/\pi_{\HMT})^2$.
& El producto $G_a\hbar_a$ conserva la normalización areal.
La construcción longitudinal y radial, sus reglas y la unicidad del
coeficiente se demuestran en \ref{g26:sec:rigidez-radial-es};
la realización circular se conserva en \ref{sec:gravity} y la confluencia térmica en
\ref{sec:confluencia}.\\
\end{longtable}
\endgroup

La comparación cuantitativa conserva el significado propio de cada
referencia. En Planck, la tabla~\ref{tab:ii-planck-comparacion} compara
coordenadas dimensionales con una constante definitoria del SI; los
residuos relativos de las orientaciones anterior y posterior son,
respectivamente, aproximadamente $-1.82\times10^{-15}$ y
$-6.93\times10^{-10}$. En Barbero--Immirzi, la
tabla~\ref{tab:ii-barbero-comparacion} compara dos construcciones:
la evaluación interna y la raíz de la ecuación de Ghosh--Mitra, cuyo
valor aproximado es $0.2740668582328300$. La diferencia relativa es
aproximadamente $-2.16\times10^{-4}$; el objeto comparado es un
coeficiente teórico.

En CKM, la tabla~\ref{tab:ii-ckm-comparacion} reúne los senos de los
tres ángulos de mezcla, la fase en radianes y el invariante de Jarlskog.
Los ángulos y la matriz de módulos se explicitan en el desarrollo
contiguo.
Las desviaciones marginales normalizadas allí definidas satisfacen
$|z|<0.029$ frente a las filas PDG utilizadas. Cada residuo utiliza
la incertidumbre marginal indicada. Su lectura es individual;
una evaluación conjunta requiere además la matriz de covarianza
de esos observables.
Para Boltzmann y gravitación, la tabla~\ref{tab:ii-kb-g-cartas}
expone el transductor, la composición longitudinal y radial, las coordenatizaciones dimensionales
y la referencia externa por separado. Esta organización permite
identificar exactamente qué cantidad produce cada construcción y
qué cantidad evalúa cada contraste.

El cuadro~\ref{tab:ii-conclusiones-contraste} reúne los resultados de
esas comparaciones junto a las definiciones estructurales anteriores.
\begin{table}[!htbp]
\centering\small
\setlength{\tabcolsep}{4pt}
\renewcommand{\arraystretch}{1.12}
\begin{tabularx}{\linewidth}{@{}>{\raggedright\arraybackslash}p{.16\linewidth}>{\raggedright\arraybackslash}X>{\raggedright\arraybackslash}X@{}}
\toprule
Magnitud & Resultado en la realización especificada & Referencia y contraste\\
\midrule
Planck, anterior al retorno
& \(\hbar_{\rm pre}=1.0545718176461544696\ldots\)
  \(\times10^{-34}\,\mathrm{J\,s}\).
& SI: \(\hbar_{\rm SI}=h_{\rm SI}/(2\pi)\), exacto.
  Su coordenada es \(1.0545718176461563912\ldots\)
  \(\times10^{-34}\,\mathrm{J\,s}\).
  Residuo relativo \(-1.822221\times10^{-15}\).\\[4pt]
Planck, posterior al retorno
& \(\hbar_{\rm ret}=1.0545718169150390611\ldots\)
  \(\times10^{-34}\,\mathrm{J\,s}\).
& Misma referencia SI.
  Residuo relativo \(-6.932836\times10^{-10}\).
  Para \(h_a=2\pi\hbar_a\), los residuos son idénticos.\\[4pt]
Barbero--Immirzi
& \(\gamma_{\HMT}=0.2740076726944592747\ldots\).
& Ghosh--Mitra: aproximadamente \(0.2740668582328300\).
  Diferencia relativa \(-2.159529202\times10^{-4}\)
  entre construcciones teóricas.\\[4pt]
Mezcla CKM
& \(\delta_{\rm CP}=65.676173123554^\circ\),
  \(J=3.1189723326632974\times10^{-5}\).
  Los tres ángulos de \eqref{eq:ii-ckm-decimales} son
  \(\theta_{12}=13.003313589509^\circ\),
  \(\theta_{23}=2.398056157332^\circ\) y
  \(\theta_{13}=0.213977382244^\circ\).
& PDG 2025: las cinco comparaciones de
  \(s_{12},s_{23},s_{13},\delta,J\) satisfacen \(|z_x|<0.029\)
  con las incertidumbres marginales de la
  tabla~\ref{tab:ii-ckm-comparacion}.\\[4pt]
Boltzmann
& \(B_*=1.380649\times10^{-23}\,u_E/u_\Theta\),
  con las bases y la regla termométrica de
  \eqref{eq:ii-b-evaluacion-comparada}.
& SI: \(1.380649\times10^{-23}\,\mathrm{J/K}\), exacto.
  Coincidencia del numeral bajo el transporte dimensional declarado.\\[4pt]
Longitud \(L\)
& \(L=1.6162549826251263301\ldots\)
  \(\times10^{-35}\,\mathrm m\), para \(L_*=1\,\mathrm m\).
& Composición longitudinal y radial:
  \(G_{\rm ret}=c_{\rm int}^3L^2/\hbar_{\rm ret}\),
  con las mismas bases de velocidad y acción.\\[4pt]
Gravitación
& \(G_{\rm ret}=6.6742996594140227\ldots\)
  \(\times10^{-11}\,\mathrm{m^3\,kg^{-1}\,s^{-2}}\).
& CODATA 2022: \(6.67430(15)\times10^{-11}\) en las mismas unidades.
  Diferencia de aproximadamente \(-0.00227057\) incertidumbres estándar.\\
\bottomrule
\end{tabularx}
\caption{Síntesis del contraste cuantitativo. Las evaluaciones, las bases
y las referencias son las de las tablas~\ref{tab:ii-planck-comparacion}--\ref{tab:ii-kb-g-cartas}.
Los residuos relativos de las constantes de acción, la comparación
teórica de Barbero--Immirzi y los residuos normalizados de CKM y gravitación
conservan sus respectivos significados.}
\label{tab:ii-conclusiones-contraste}
\par\medskip
\centering\small
\setlength{\tabcolsep}{4pt}
\renewcommand{\arraystretch}{1.18}
\begin{tabular}{@{}lrrr@{}}
\toprule
Observable & Evaluación HMT & PDG 2025 & \(z_x\)\\
\midrule
\(s_{12}\) & \(0.225007404757\) & \(0.22501\pm0.00068\) & \(-0.003817\)\\
\(s_{23}\) & \(0.041841757010\) & \(0.04183^{+0.00079}_{-0.00069}\) & \(0.014882\)\\
\(s_{13}\) & \(0.003734601164\) & \(0.003732^{+0.000090}_{-0.000085}\) & \(0.028902\)\\
\(\delta\) (rad) & \(1.146265461116\) & \(1.147\pm0.026\) & \(-0.028251\)\\
\(10^5J\) & \(3.118972333\) & \(3.12^{+0.13}_{-0.12}\) & \(-0.008564\)\\
\bottomrule
\end{tabular}
\caption{Valores CKM y comparación marginal reunidos en las conclusiones.
Los tres ángulos de la tabla~\ref{tab:ii-conclusiones-contraste} se contrastan
mediante \(s_{ij}=\sin\theta_{ij}\), junto con la fase en radianes y el
invariante \(J\). Se conservan las cifras, incertidumbres y residuos de la
tabla~\ref{tab:ii-ckm-comparacion}, correspondientes al ajuste de tres
generaciones PDG 2025. La última fila expresa conjuntamente valor e
incertidumbre en unidades de \(10^{-5}\). Los residuos son comparaciones
marginales de observables correlacionados.}
\label{tab:ii-conclusiones-ckm}
\end{table}


\Needspace{10\baselineskip}
\subsection{Conservación de la escala areal y recuperación de la información de frontera}

La relación entre información y gravitación se concentra en dos mapas
complementarios. El primero conserva la escala areal mediante
$G_a\hbar_a$ y organiza sus ocupaciones mediante
$\gamma_{\HMT}a_0$. El segundo conserva la resolución sectorial mediante
la observación conjunta del área y de la incidencia ponderada. La inversión
\eqref{eq:recover} recupera las ocupaciones desde el área y el
observable ponderado; las ecuaciones
\eqref{eq:ii-respuesta-entropica} y \eqref{eq:ii-respuesta-areal}
vinculan sus variaciones con las fluctuaciones del estado.

De aquí procede la lectura física central del artículo: una misma
configuración admite una observación geométrica y una observación
informacional relacionadas por una estructura recuperable. El cambio
de orientación puede conservar la entropía del estado reducido y
transformar el área, como establecen
\eqref{eq:ii-inversion-entropia} y \eqref{eq:ii-area-complementaria}.
La conservación de información se refiere al estado y a su transporte;
la entropía de entrelazamiento se refiere además a la bipartición y a
la reducción. La identidad de horizonte de
\ref{sec:confluencia} utiliza esas funciones en la realización
geométrica y térmica especificada.

Las identidades de memoria permiten precisar también el alcance de una
reducción de variables. La igualdad de observaciones de
\eqref{mem:schur-lector} conserva las contribuciones del sector eliminado
mediante el operador efectivo, la fuente y el lector transformados.
La entropía del estado reducido, en cambio, se evalúa con la bipartición
y la medida especificadas en \ref{sec:ii-ecuacion-informacion}.
Ambas construcciones explican funciones diferentes de la memoria:
la primera conserva una lectura del transporte; la segunda cuantifica
la información accesible al subsistema. Su articulación mantiene
explícito qué estructura se transporta y qué reducción define el
observable físico.

El avance explicativo se expresa, por tanto, mediante un orden preciso:
estructura generada, aplicación de lectura, valor obtenido y función
física. La igualdad de dos valores se estudia junto con la igualdad
o diferencia de sus estructuras; los mapas de recuperación determinan
cuál de esas diferencias conserva significado observable.
% END THERMAL_R03_SYNC_13B


Las comparaciones cuantitativas mantienen la distinción entre una
constante definitoria de unidades, un coeficiente teórico y observables
de mezcla con incertidumbres. Examinan las coordenadas producidas en
las coordenatizaciones declaradas, mientras las identidades estructurales determinan
qué se conserva al transportarlas: el cociente de retorno de la acción,
la orientación del funcional de Barbero--Immirzi, la métrica y la fase
de la mezcla, la transducción térmica y el producto gravitatorio.

La contribución del artículo es esta determinación articulada de las
constantes y de sus relaciones físicas. La composición APP--TRIT--TPK
construye las regiones numéricas y la memoria dodecafásica antes de sus
evaluaciones; los lectores posteriores producen las secciones de acción,
las coordenadas angulares y los observables de frontera. En la realización
de horizonte especificada, \(T_HS_H=E_\Lambda\) y la celda de radio
\(\ell_P\) reúne energía de decisión y media acción. La acción fija
energía por duración; Barbero--Immirzi resuelve la incidencia areal;
Boltzmann relaciona energía e información; y la gravitación conserva
la unidad geométrica \(G_a\hbar_a/c_{\rm int}^3\).
Estas relaciones expresan el significado de las constantes dentro de
una estructura común y conservan las condiciones precisas bajo las que
cada una queda determinada.
% END THERMAL_R03_SYNC_13
